\documentclass[../main]{subfiles}
\begin{document}
\chapter{Estudio de tiempos}
\img{images/11/estudio_tiempos.jpg}{Estudio de tiempos.}
\info{Estudio de tiempos}{Un Estudio de tiempos y movimientos es una
  técnica de eficiencia en el negocio que combina el trabajo de Estudio
  de Tiempos realizado por Frederick Winslow Taylor junto con el
  trabajo de Estudio de Movimientos de Frank y Lillian Gilbreth. Es un
trabajo mayoritariamente de la Administración científica.}

\subsection{Procedimiento directo de estudio de tiempos}
A continuación se muestra el procedimiento desarrollado por Mikell
Groover para un estudio de tiempos directo
\begin{enumerate}

  \item  Definir y documentar el método estándar.
  \item  Dividir la tarea en elementos de trabajo. Estos primeros dos
    pasos son prioridad para elegir el ritmo oportuno. Familiarizan
    al analista con la tarea y le permiten intentar mejorar el
    procedimiento de trabajo antes de definir el tiempo estándar.
  \item  Cronometrar los elementos de trabajo para obtener el tiempo
    observado para la tarea.
    Evaluar el ritmo del trabajador relativo al desempeño estándar
    (clasificación del desempeño), para determinar el tiempo normal.
  \item Tome nota de que los pasos 3 y 4 se completan de manera
    simultánea. Durante estos pasos, diferentes ciclos de trabajo son
    cronometrados y el desempeño de cada ciclo es calificado por
    separado. Finalmente, los datos obtenidos en estos pasos son
    promediados para generar el tiempo normalizado.
  \item  Aplicar un margen de error al tiempo normal para calcular el
    tiempo estándar. Los márgenes de factores necesarios en el
    trabajo son añadidos para calcular el tiempo estándar de la tarea.

\end{enumerate}
\subsection{El estudio de tiempos}

El estudio de tiempos, también conocido como el método clásico con
cronómetro, fue propuesto por Frederick Taylor en 1881. Aunque se han
desarrollado otras metodologías para medir el trabajo, el método
clásico con cronómetro sigue siendo el más utilizado. Este estudio
consiste en medir el tiempo que un trabajador dedica a realizar una
tarea determinada, con el objetivo de establecer un tiempo estándar.

Es importante tener en cuenta que el Estudio de Tiempos y la Medición
del trabajo no son sinónimos. La Medición del trabajo se refiere a la
aplicación de técnicas para determinar el tiempo que un trabajador
calificado dedica a realizar una tarea específica de acuerdo con una
norma de ejecución preestablecida. El Estudio de Tiempos es una de
estas técnicas.

\info{La medición del trabajo según la OIT (Organización
Internacional del Trabajo)}{La Medición del trabajo es la aplicación
  de técnicas para determinar el tiempo que invierte un trabajador
  calificado en llevar a cabo una tarea definida efectuándola según una
norma de ejecución preestablecida.}

\subsubsection{¿Cuál es el objetivo de la Medición del trabajo?}

Como se puede ver en el módulo de Estudio del Trabajo, el tiempo de
trabajo puede aumentar debido a un mal diseño del producto, un mal
funcionamiento del proceso o debido a tiempo improductivo atribuible
a la gestión o a los trabajadores. El Estudio de Métodos es la
 técnica por excelencia para minimizar la cantidad de trabajo de tipo
operativo, eliminar movimientos innecesarios y reemplazar métodos. La
Medición del trabajo, por su parte, sirve para investigar, minimizar
y eliminar el tiempo ocioso del proceso.

Una función adicional de la Medición del Trabajo es la fijación de
tiempos estándar (tiempos tipo) de ejecución, por lo que es una
herramienta complementaria en la Ingeniería de Métodos, sobre todo en
las fases de definición e implantación. Además de ser una herramienta
invaluable para el costeo de las operaciones.

\info{Enfoque en la medición}{Es importante tener un enfoque
  sistémico al considerar tanto técnicas de estudio de métodos como
  de medición del trabajo. Es decir, considerar el rendimiento de las
  operaciones en conjunto, ya que la organización es un sistema, un
  todo indivisible.
  Es fundamental evitar centrarse en mejorar el rendimiento local sin
  tener en cuenta el objetivo del sistema.}

Las personas, parte fundamental de las organizaciones, no suelen
rechazar el cambio en sí mismo, sino que rechazan el cambio del cual
no perciben un beneficio. Históricamente, el estudio de tiempos fue
utilizado como una técnica de control de lo que se consideraban
causas de ineficiencia atribuibles a los trabajadores, lo que le dio
mala reputación, especialmente en los círculos sindicales.

Comprender y transmitir que los resultados de un estudio de tiempos
pueden ser potencialmente utilizados para la gestión de operaciones,
la administración del flujo, el equilibrio de cargas, entre otros; y
que, en consecuencia, pueden beneficiar a la organización y, por
ende, a los trabajadores, es fundamental en este ejercicio, ya que
implica partir desde consideraciones humanas.

En el desarrollo de un profesional habrá muchas ocasiones en las que
requerirá alguna técnica de medición del trabajo. En el proceso de
fijación de tiempos estándar es posible que sea necesario utilizar la
medición para:

\begin{enumerate}
  \item Comparar la eficacia de varios métodos, y elegir el que
    requiera menor tiempo de ejecución en igualdad de condiciones.
  \item Repartir el trabajo dentro de los equipos con la ayuda de
    diagramas de actividades múltiples, con el objetivo de equilibrar
    los procesos.
  \item Determinar el número de máquinas que puede atender un operario.
\end{enumerate}

Una vez se ha determinado el tiempo estándar (tipo), este puede utilizarse para:
\begin{enumerate}
  \item Obtener la información de base para el programa de producción.
  \item Obtener información para basar cotizaciones, precios de venta
    y plazos de entrega.
  \item Fijar normas sobre el uso de la maquinaria y la mano de obra.
  \item Obtener información que permita controlar los costos de la
    mano de obra (incluso establecer planes de incentivos) y mantener
    costos estándar.

\end{enumerate}

\subsection{Técnicas de Medición del trabajo}

Cuando mencionábamos que el término Medición del Trabajo no era
equivalente al término Estudio de Tiempos, nos referimos a que el
Estudio de Tiempos es tan solo una de las técnicas contenidas en el
conjunto «Medición». Las principales técnicas que se emplean en la
medición del trabajo son:

\begin{enumerate}
  \item Muestreo del Trabajo
  \item Estimación Estructurada
  \item Estudio de Tiempos
  \item Normas de Tiempo Predeterminadas
  \item Datos Tipo
\end{enumerate}

\subsection{Procedimiento básico sistemático para realizar una
Medición del Trabajo}
Las etapas necesarias para efectuar sistemáticamente la medición
del trabajo son:

\begin{enumerate}
  \item Seleccionar: El trabajo que va a ser objeto de estudio.
  \item Registrar: Todos los datos relativos a las circunstancias en
    que se realiza el trabajo, a los métodos y a los elementos de
    actividad que suponen.
  \item Examinar: Los datos registrados y el detalle de los elementos
    con sentido crítico para verificar si se utilizan los métodos y
    movimientos más eficaces, y separar los elementos improductivos o
    extraños de los productivos.
  \item Medir: La cantidad de trabajo de cada elemento, expresándola
    en tiempo, mediante la técnica más apropiada de medición del trabajo.
  \item Compilar: El tiempo estándar de la operación previendo, en
    caso de estudio de tiempos con cronómetro, suplementos para
    breves descansos, necesidades personales, etc.
  \item Definir: Con precisión la serie de actividades y el método de
    operación a los que corresponde el tiempo computado y notificar
    que ese será el tiempo estándar para las actividades y métodos
    especificados.

\end{enumerate}

\info{¿Qué es el Estudio de tiempos?}{El Estudio de Tiempos es una
  técnica de medición del trabajo empleada para registrar los tiempos y
  ritmos de trabajo correspondientes a los elementos de una tarea
  definida, efectuada en condiciones determinadas y para analizar los
  datos a fin de averiguar el tiempo requerido para efectuar la tarea
según una norma de ejecución preestablecida.}

Es innegable que dentro de las técnicas que se emplean en la medición
del trabajo la más importante es el Estudio de Tiempos, o por lo
menos es la que más nos permite confrontar la realidad de los
sistemas productivos sujetos a medición.

\actividad{Investigar y responder}
\begin{enumerate}
  \item ¿Qué es el estudio de tiempos?
  \item ¿Cuáles son los pasos para realizar un estudio de tiempos?
  \item Resumir el procedimiento directo de estudio de tiempos que
    indica el texto.
  \item ¿Por qué es importante el estudio de tiempos?
  \item ¿Qué se busca lograr al realizar un estudio de tiempos?
  \item ¿Cómo se mide el tiempo en un estudio de tiempos?
  \item ¿Qué beneficios puede aportar un estudio de tiempos a una empresa?
  \item ¿Qué herramientas simples se pueden utilizar para llevar a
    cabo un estudio de tiempos?
  \item ¿Cuál es el propósito de establecer tiempos estándar?
  \item ¿Qué factores pueden afectar la precisión de un estudio de tiempos?
  \item ¿Cómo se pueden mejorar los procesos a través del estudio de tiempos?
  \item ¿Cuál es la importancia de la eficiencia en las operaciones
    de una empresa?
  \item Resumir el procedimiento básico sistemático para realizar una
    Medición de trabajo.
\end{enumerate}
\end{document}